\section{One iteration}
\label{sec:iteration}

The state at the top of an outer iteration is a working set $\Active$ of size
$k$, the factors $(J,R)$ satisfying~\eqref{eq:inv1}--\eqref{eq:inv2}, the
subproblem solution $x$ and multipliers $u$ of Proposition~\ref{prop:sub}, and
the running objective $f(x)$. The invariant maintained throughout is
\begin{equation}
  \text{$(x,u)$ solves the subproblem~\eqref{eq:sub}, and}
  \quad u_i \ge 0 \ \text{ for every inequality } i \in \Active,
  \label{eq:invariant}
\end{equation}
which by the discussion after Proposition~\ref{prop:sub} is stationarity,
complementarity and dual feasibility together. Only primal feasibility is
missing, and the iteration exists to obtain it.

\subsection{Choosing the entering constraint}

Compute the slacks $s_i = c_i\T x - b_i$ for every $i$. If $s_i \ge 0$ for all
inequalities and $s_i = 0$ for all equalities, then~\eqref{eq:pfeas} holds,
\eqref{eq:invariant} supplies the rest, and $x$ is the solution by
Proposition~\ref{prop:sufficient}. Otherwise a violated constraint is selected as
the entering constraint. The rule is the most violated one, measured relative to
the norm of its own normal:
\begin{equation}
  v_i = \begin{cases} |s_i|, & i \le m_{\mathrm{eq}},\\ -s_i, & i > m_{\mathrm{eq}},\end{cases}
  \qquad
  i_* \in \argmax_i \; \frac{v_i}{\norm{c_i}},
  \qquad\text{terminating when } \max_i \frac{v_i}{\norm{c_i}} \le 0,
  \label{eq:select}
\end{equation}
so that an equality counts as violated in either direction while an inequality
counts only when its slack is negative. The scaling is not cosmetic.

\begin{proposition}[Selection is scale invariant]\label{prop:scale}
Replacing $(c_i, b_i)$ by $(\gamma c_i, \gamma b_i)$ for some $\gamma > 0$ leaves
the feasible set, the solution, and the choice made by~\eqref{eq:select}
unchanged.
\end{proposition}

\begin{proof}
The constraint $\gamma c_i\T x \ge \gamma b_i$ has the same solution set as
$c_i\T x \ge b_i$. Its slack scales as $s_i \mapsto \gamma s_i$ and its normal as
$\norm{c_i} \mapsto \gamma\norm{c_i}$, so the ratio in~\eqref{eq:select} is
unchanged, as are the other constraints' ratios.
\end{proof}

\begin{remark}[The invariance is exact, and covers the whole walk]
\label{rem:scaleexact}
Proposition~\ref{prop:scale} claims only that the \emph{choice} is unchanged.
Measured against the implementation over three decades of $\gamma$ either side of
unity, the statement is stronger: the solver performs an identical number of
additions and removals, so the whole trajectory is unchanged and not only its first
step, and returns a minimiser agreeing to $3.3\times10^{-16}$.
\end{remark}

Without the division, a constraint written in basis points rather than in units
would be selected first merely for having been written differently, and the
solver's trajectory, though not its answer, would depend on the caller's
choice of units. Note that the multiplier of the rescaled constraint scales as
$\lambda_i \mapsto \lambda_i/\gamma$, so the sign conditions are unaffected too.

\subsection{The two directions}

Write $n_* = c_{i_*}$ and $b_* = b_{i_*}$, and form $d$, $z$ and $r$
as in~\eqref{eq:dzr}. The following lemma is the computational core of the
method; everything in the rest of this section is a corollary of it.

\begin{lemma}[Properties of the directions]\label{lem:directions}
With $d$, $z$, $r$ as in~\eqref{eq:dzr},
\begin{enumerate}
  \item[(i)] $A\T z = 0$: moving along $z$ preserves the working-set equalities;
  \item[(ii)] $n_*\T z = \norm{d_2}^2 = z\T G z \ge 0$, with equality iff $z = 0$;
  \item[(iii)] $G z = n_* - A r$;
  \item[(iv)] $z = 0$ if and only if $n_* = A r$, i.e.\ iff $n_*$ lies in the span of
        the working-set normals, in which case $r$ holds its coordinates there.
\end{enumerate}
\end{lemma}

\begin{proof}
(i) $A\T z = A\T J_2 d_2 = (J_2\T A)\T d_2 = 0$ by~\eqref{eq:inv2}.

(ii) $n_*\T z = n_*\T J_2 d_2 = (J_2\T n_*)\T d_2 = d_2\T d_2$. For the second
equality, $z\T G z = d_2\T (J_2\T G J_2) d_2 = d_2\T d_2$ by the trailing block
of~\eqref{eq:inv1}. Since $J_2$ has full column rank, $z = 0$ iff $d_2 = 0$.

(iii) By~\eqref{eq:reduced}, $Gz = GJ_2J_2\T n_* = n_* - A(A\T G^{-1}A)^{-1}A\T
G^{-1}n_*$, and the second term is $Ar$ by~\eqref{eq:rclosed}.

(iv) If $d_2 = 0$ then $J\T n_* = (d_1,0) = [R;0]R^{-1}d_1 = (J\T A) r$, and $J$
is nonsingular, so $n_* = Ar$. Conversely $n_* = Ar$ gives
$Gz = n_* - Ar = 0$ by (iii), hence $z = 0$.
\end{proof}

Part (ii) identifies $z$ as an ascent direction for the entering constraint's
slack: the slack $n_*\T x - b_*$ grows at rate $n_*\T z > 0$ per unit step. Part
(i) says the step costs nothing in working-set feasibility. Together with the
closed forms~\eqref{eq:zclosed} and~\eqref{eq:rclosed} the reading is:
\begin{itemize}
  \item $z$ is the $G$-orthogonal projection of the unconstrained direction
        $G^{-1}n_*$ onto the working set's feasible subspace $\nullsp(A\T)$.
        Indeed $\Pi := J_2 J_2\T G$ is idempotent with range $\nullsp(A\T)$ and is
        self-adjoint for $\langle\cdot,\cdot\rangle_G$, and
        $z = \Pi\, G^{-1}n_*$.
  \item $r$ solves a least-squares problem in the $G^{-1}$ metric. By
        \eqref{eq:rclosed},
        \begin{equation}
          r = \argmin_{y \in \R^k} \; \norm{J\T A y - J\T n_*}_2
            = \argmin_{y \in \R^k} \; (Ay - n_*)\T G^{-1} (A y - n_*),
          \label{eq:lsq}
        \end{equation}
        the best approximation of the entering normal by the working-set normals.
        Its residual, mapped through $G^{-1}$, is $z$: from
        Lemma~\ref{lem:directions}(iii), $z = G^{-1}(n_* - Ar)$.
\end{itemize}

\begin{remark}[Why not call a least-squares routine]
Equation~\eqref{eq:lsq} is a genuine least-squares problem and it is tempting to
hand it to a library. The temptation should be resisted: the factor $R$ of its
normal equations is already held, so~\eqref{eq:dzr} costs one triangular solve,
where a library call would refactorise $J\T A$ from scratch at $O(nk^2)$ every
iteration. The same remark applies to~\eqref{eq:subsol}: the closed form is the
specification, not the implementation.
\end{remark}

\subsection{The affine path}

Fix the directions and consider moving along them. Let $u_*$ denote the entering
constraint's own multiplier, which starts at $0$.

\begin{proposition}[Dual feasibility along the path]\label{prop:path}
For $t \in \R$ define
\begin{equation}
  x(t) = x + t z, \qquad u(t) = u - t r, \qquad u_*(t) = u_* + t .
  \label{eq:path}
\end{equation}
Then for all $t$: $A\T x(t) = b_\Active$, and stationarity holds for the working
set augmented by the entering constraint,
\begin{equation}
  G x(t) - a = A\, u(t) + u_*(t)\, n_* .
  \label{eq:pathstat}
\end{equation}
\end{proposition}

\begin{proof}
The first claim is Lemma~\ref{lem:directions}(i). For the second, stationarity at
$t = 0$ reads $Gx - a = Au + u_* n_*$, and by Lemma~\ref{lem:directions}(iii),
\[
  Gx(t) - a = (Gx - a) + t\,Gz
  = A u + u_* n_* + t(n_* - Ar)
  = A(u - tr) + (u_* + t) n_* . \qedhere
\]
\end{proof}

So the entire path is stationary, with the entering multiplier rising at unit rate
and the working-set multipliers falling linearly along $-r$. Complementarity holds
for every constraint except the entering one, whose slack is still nonzero while
its multiplier is positive; that single violation is what the step is taken to
remove. Two limits on $t$ follow immediately.

\paragraph{The dual limit $t_1$.} Dual feasibility requires $u_i(t) \ge 0$ for
every inequality $i \in \Active$. Since $u_i(t) = u_i - t r_i$ and $u_i \ge 0$, the
binding indices are those with $r_i > 0$, and
\begin{equation}
  t_1 = \min\Bigl\{\, \frac{u_i}{r_i} \;:\; i \in \Active,\ i > m_{\mathrm{eq}},\
  r_i > 0 \,\Bigr\},
  \qquad
  i_{\mathrm{del}} = \text{an index attaining it},
  \label{eq:t1}
\end{equation}
with $t_1 = +\infty$ when the set is empty. Equality constraints are excluded:
their multipliers are unrestricted in sign, so no step in $t$ can make them
infeasible. This is the classical ratio test, and $t_1$ is the largest step that
preserves~\eqref{eq:invariant}.

\paragraph{The primal limit $t_2$.} The entering constraint's slack is
$n_*\T x(t) - b_* = s_{i_*} + t\, n_*\T z$ with $s_{i_*} < 0$, so it reaches zero
at
\begin{equation}
  t_2 = \frac{|s_{i_*}|}{n_*\T z} = \frac{|s_{i_*}|}{\norm{d_2}^2},
  \label{eq:t2}
\end{equation}
finite and positive whenever $z \neq 0$, and $t_2 = +\infty$ when $z = 0$ by
Lemma~\ref{lem:directions}(ii), in which case the primal iterate cannot move at all
and only the multipliers change.

\paragraph{The step.} Take $t = \min(t_1,t_2)$. If $t_2 \le t_1$ the step is
\emph{full}: the entering constraint now holds with equality, so it joins the
working set with multiplier $u_* + t_2$ and the factorisation is updated
(Section~\ref{sec:updates}). The new state satisfies~\eqref{eq:invariant}:
stationarity and complementarity by Proposition~\ref{prop:path} with the slack now
zero, and the sign conditions because $t_2 \le t_1$. If $t_1 < t_2$ the step is
\emph{partial}: the multiplier of $i_{\mathrm{del}}$ reaches zero first, that
constraint leaves the working set, the factorisation is downdated, and $z$ and $r$
are recomputed, both having changed with $J$ and $R$, keeping the same
entering constraint and its accumulated $u_*$. That is the inner loop. If both
limits are infinite the problem is infeasible; see
Proposition~\ref{prop:farkas}.

\begin{remark}[Equalities violated from above]\label{rem:reverse}
An equality constraint may be violated in either direction. When $s_{i_*} > 0$
the slack must be \emph{decreased}, so the step is taken along $-z$: formally
$t$ ranges over the negative reals, $t_2 = -|s_{i_*}|/n_*\T z$, the ratio test in
\eqref{eq:t1} runs against $-r$ rather than $r$, and the entering multiplier
accumulates negatively. Nothing else changes; in particular
Proposition~\ref{prop:path} was stated for all $t \in \R$ precisely so that this
case needs no separate treatment, and Proposition~\ref{prop:increment} below
covers it too.
\end{remark}

\subsection{The objective, in closed form}

The running objective is not re-evaluated from $x$; it is advanced by an exact
increment.

\begin{proposition}[Objective increment]\label{prop:increment}
Along the path~\eqref{eq:path},
\begin{equation}
  f(x(t)) - f(x) = t\Bigl(\frac{t}{2} + u_*\Bigr)\, n_*\T z ,
  \label{eq:increment}
\end{equation}
which is strictly positive whenever $z \neq 0$ and $t \neq 0$ and $t$, $u_*$ have
the same sign (in particular whenever $t > 0$ and $u_* \ge 0$).
\end{proposition}

\begin{proof}
Expanding the quadratic,
$f(x + tz) = f(x) + t\, z\T(Gx - a) + \tfrac12 t^2\, z\T G z$. For the linear
term, stationarity at $t = 0$ and Lemma~\ref{lem:directions}(i) give
$z\T(Gx - a) = z\T A u + u_*\, z\T n_* = u_*\, n_*\T z$. For the quadratic term,
Lemma~\ref{lem:directions}(ii) gives $z\T G z = n_*\T z$. Hence
$f(x+tz) - f(x) = t u_* n_*\T z + \tfrac12 t^2 n_*\T z$, which
is~\eqref{eq:increment}. Positivity: $n_*\T z > 0$ by
Lemma~\ref{lem:directions}(ii), and $t(t/2 + u_*) > 0$ when $t$ and $u_*$ share a
sign and $t \neq 0$.
\end{proof}

Three things are worth noting about~\eqref{eq:increment}. It is exact, so the
objective can be carried as a running total at the cost of one multiplication per
step instead of an $O(n^2)$ re-evaluation of the quadratic. It is expressed in
quantities the iteration already holds, $n_*\T z$ being needed for $t_2$ anyway. And it handles the reversed step of Remark~\ref{rem:reverse} without a
special case: there $t < 0$ and $u_*$ accumulates negatively, so
$t/2 + u_* < 0$ and the product with $t$ is again positive. The value therefore
increases on every nonzero step, in both directions, which is the monotonicity
Section~\ref{sec:termination} needs.

\begin{algorithm}[ht]
\caption{The Goldfarb--Idnani dual active-set method}\label{alg:gi}
\begin{algorithmic}[1]
\Require $G \succ 0$, $a$, $C$, $b$, $m_{\mathrm{eq}}$
\State $J \leftarrow \mathrm{chol}(G)^{-1}$; \ $x \leftarrow G^{-1}a$; \
       $f \leftarrow -\tfrac12 a\T x$; \ $\Active \leftarrow \emptyset$; \
       $k \leftarrow 0$; \ $u \leftarrow ()$
\Loop
  \State $s \leftarrow C\T x - b$, with $s_i \leftarrow 0$ for $i \in \Active$;
         choose $i_*$ by~\eqref{eq:select}
  \State \textbf{if} none \textbf{then} \Return $(x,f,u)$
         \Comment{optimal, by Prop.~\ref{prop:sufficient}}
  \State $n_* \leftarrow c_{i_*}$; \ $u_* \leftarrow 0$
  \Loop
    \State $d \leftarrow J\T n_*$; \ $z \leftarrow J_2 d_2$; \
           $r \leftarrow R^{-1} d_1$
    \State $t_1, i_{\mathrm{del}} \leftarrow$ \eqref{eq:t1}; \
           $t_2 \leftarrow$ \eqref{eq:t2};
           \textbf{if} both infinite \textbf{then stop}: infeasible
    \State $t \leftarrow \min(t_1,t_2)$; \ $x \leftarrow x + tz$; \
           $f \leftarrow f + t(t/2 + u_*)\,n_*\T z$; \
           $u \leftarrow u - tr$; \ $u_* \leftarrow u_* + t$
    \If{$t_2 \le t_1$}
      \State $\Active \leftarrow \Active \cup \{i_*\}$, $u_{k+1} \leftarrow u_*$,
             $k \leftarrow k+1$; insert (\S\ref{ssec:insert}); \textbf{break}
    \Else
      \State $\Active \leftarrow \Active \setminus \{i_{\mathrm{del}}\}$,
             $k \leftarrow k-1$; delete (\S\ref{ssec:delete});
             $s_{i_*} \leftarrow n_*\T x - b_*$
    \EndIf
  \EndLoop
\EndLoop
\end{algorithmic}
\end{algorithm}
